% C1b/Q1: source-preserved standalone support lemma extracted from V27.
\section{Local irreducible supports in binary/unary partial structures}
\label{sec:binary-unary-irreducible-local-cover}

\begin{lemma}[Binary/unary irreducible substructures are covered pairwise by local closed supports]
\label{lem:typed-binary-irreducible-local-cover}
Let \(C\) be a finite
PARTIAL relational
structure with total
named UNARY functions
and relational symbols
of arity AT MOST TWO,
closed under all the
named functions.
For every vertex \(v\in C\)
form the closed support
\(\cl_C\{v\}\).
For every ACTUALLY
DISPLAYED binary relation
tuple on two distinct
vertices \(u,v\in C\),
also form the closed
support \(\cl_C\{u,v\}\).
Let \(\mathscr H(C)\)
be this finite family
of induced functionally
closed substructures.

Then for EVERY nonempty
functionally closed
IRREDUCIBLE induced
substructure \(D\subseteq C\),
and EVERY pair of distinct
vertices \(x,y\in D\),
SOME member of
\(\mathscr H(D)\)
contains BOTH \(x,y\).

Consequently, if
\(f:C\to H\) is a
sort- and
unary-function-preserving
homomorphism to a
COMPLETE target, and
\(f|_A\) is an induced
EMBEDDING for every
\(A\in\mathscr H(C)\),
then \(f\) is a
HOMOMORPHISM-EMBEDDING
on all of \(C\):
its restriction to EVERY
source irreducible
substructure is an
embedding.

This conclusion is
specific to the binary
relational/unary-functional
signature. For a higher
arity relation, pairwise
reflection does NOT
by itself reflect a
missing tuple of
three or more roots.
\end{lemma}
\begin{proof}
All members of
\(\mathscr H(D)\)
are functionally closed:
the closure of a finite
set under UNARY
functions is the
union of the
one-generated closures
of its elements.
Every original
displayed unary-function
equation is contained
in a one-generated
member of
\(\mathscr H(D)\),
and every displayed
binary relation tuple
is contained in the
closed support of its
two endpoints.

Suppose for contradiction
that distinct \(x,y\in D\)
lie together in NO
member of
\(\mathscr H(D)\).
Define two subsets
\[
 \begin{aligned}
 A_x&=\bigcup\{U\in\mathscr H(D):
                              y\notin U\},\\
 A_y&=\bigcup\{U\in\mathscr H(D):
                              x\notin U\}.
 \end{aligned}
\]
Each is functionally
closed, because it is
a union of sets closed
under unary functions.
We have \(x\in A_x
\setminus A_y\)
and \(y\in A_y
\setminus A_x\):
the one-generated
support of \(x\)
avoids \(y\), and
that of \(y\) avoids
\(x\), whereas all
sets in the union
defining \(A_y\)
avoid \(x\), and
vice versa.

Every vertex \(z\in D\)
belongs to its own
one-generated support
\(\cl_D\{z\}\).
That support cannot
contain BOTH \(x,y\)
by the assumption,
so it occurs in at
least one of the two
unions. Hence
\(D=A_x\cup A_y\).

Every displayed binary
relation tuple of D
lies in its own
closed relation support
\(U\in\mathscr H(D)\).
Again \(U\) cannot
contain both \(x,y\),
so the entire tuple
lies in \(A_x\)
or in \(A_y\).
Unary functions likewise
stay inside each
functionally closed part.
Consequently the
PARTIAL diagram of D
is exactly the free
amalgam of its two
induced closed
substructures on
\(A_x,A_y\) over
their induced closed
intersection.
Both sides have
a vertex outside
the overlap, so
this is a NONTRIVIAL
free decomposition.
This contradicts
irreducibility of D.

For the last assertion,
fix an arbitrary
irreducible induced
substructure \(D\subseteq C\).
Each member of
\(\mathscr H(D)\)
is an induced
substructure of
a corresponding
member of
\(\mathscr H(C)\):
unary closure inside
D equals closure
inside C because
D is functionally
closed. Therefore
f is an induced
embedding on every
member of
\(\mathscr H(D)\).
For any distinct
\(x,y\in D\), the
first part supplies
one such common
closed support.
Its embedding property
proves \(f(x)\ne f(y)\).
Thus \(f|_D\) is
injective.

The homomorphism f
already preserves
all unary functions
and displayed relations.
To reflect any
ACTUAL binary relation
\(R(f(x),f(y))\)
of the target for
vertices \(x,y\in D\),
choose a local
closed support
containing both;
f is an induced
embedding there,
so \(R(x,y)\)
was displayed in
that support and
hence in D.
For \(x=y\) use
the one-generated
support of x instead.
This proves reflection
of every binary and
unary relation, and
therefore \(f|_D\)
is an embedding.
As D was arbitrary,
f is a
homomorphism-embedding.
\end{proof}


\begin{lemma}[HN irreducibility forces completed unary-generated diagrams]
\label{lem:hn-irreducible-forces-unary-admissible}
Let $L$ be a finite many-sorted language with arbitrary
relational symbols and only UNARY function symbols.
In this lemma, function symbols of a finite PARTIAL
$L$-structure $C$ are permitted to have proper domains,
as in the HN 2019 finite-model convention.
An induced $L$-substructure means an EMBEDDED one:
its inclusion reflects not only relations but also
the DOMAINS of all partially defined functions.
Thus a substructure containing $x$ necessarily
contains the value of every unary function which
is already defined at $x$ in $C$.

Let $\mathcal K$ be any class of finite COMPLETED
$L$-structures whose named unary functions are
TOTAL on their intended input sorts.
Assume EVERY irreducible induced $L$-substructure
of $C$ is an ACTUAL member of $\mathcal K$.
Then:
\begin{enumerate}[label=(\roman*)]
\item For every $x\in C$, the substructure
$U_x=\cl_C\{x\}$ generated by all iterated
ALREADY DEFINED unary values is IRREDUCIBLE,
and belongs to $\mathcal K$.
\item Consequently every intended unary function
is already TOTAL in $C$ on its named input sort,
and every actual unary term identity valid in
$\mathcal K$ holds in $C$ pointwise.
\item If all completed one-generated $\mathcal K$
members have size at most $m$, then
$|U_x|\le m$ for EVERY $x\in C$, and
$|\cl_C(B)|\le m|B|$ for each finite
generator set $B\subseteq C$.
\end{enumerate}
In particular the HN local-finiteness requirement
that every irreducible substructure of $C$
already belong to the completed target age
implies the unary-admissibility condition needed
for the RESTRICTED hidden bounded-obstruction
theorem, even though the ambient HN language
permits PARTIAL unary functions.
\end{lemma}
\begin{proof}
An HN induced substructure is closed under
EVERY unary value already defined on its points,
since embedding requires equality of FUNCTION
DOMAINS on all input tuples from that substructure.
Therefore the finite intersection of all such
substructures containing a given $x$ is the
substructure $U_x$ consisting precisely of the
iterated defined unary values of $x$.

Suppose $U_x$ were the free amalgam of two
proper induced substructures over their
common induced intersection. One of these
two factors contains $x$. As that factor
is an ACTUAL substructure, it is closed
under all defined unary functions, and
hence contains ALL of $U_x$. Thus it is
not proper, a contradiction. So $U_x$
is irreducible. The hypothesis yields
$U_x\in\mathcal K$.

Because the function-domain interpretations
in each such completed $U_x$ are TOTAL on
their intended source sorts, every applicable
function value at $x$ is already defined in
$U_x$ and hence in $C$. This holds for each
$x$, proving the totality assertion.
The source unary term evaluation on any one
$x$ takes place entirely inside the ACTUAL
completed $U_x$, so every universal unary
identity valid in $\mathcal K$ holds in $C$.

Finally $U_x\in\mathcal K$ gives $|U_x|\le m$.
With unary functions the closure of any $B$
is the union of the one-generated closures
$U_x$, $x\in B$. This yields the asserted
finite bound, without assuming the original
PARTIAL source $C$ was complete or completable.
\end{proof}
